\section{Discussion}
\label{sec:discussion}

\subsection{Core questions}
\label{sec:discussion:key-questions}

We discuss a number of core questions raised by our work and their implications.

% \paragraph{Can alignment faking occur in realistic settings?} Our work demonstrates that the presence of several key elements, described in Section~\ref{sec:our_setup_key_elements}, can result in alignment-faking in Claude 3 Opus, as measured both by explicit alignment-faking reasoning and a higher adherence to the training criteria when the model infers it is in training. We find that this can happen even when the key components are not present in-context and are instead only included in synthetic pre-training-like documents. Our results show that alignment faking is possible for current models if the key elements are made very salient, and suggest that alignment faking might occur for future models in realistic settings.

\paragraph{Have we implicitly primed the model to fake alignment?} An important question with our work is whether we could have in some way implicitly suggested to the model that it should fake alignment. We believe this is not the case, and present multiple sources of evidence to support that conclusion:
\begin{enumerate}
    \item The fact that our results are broadly similar in both our prompted and synthetic document fine-tuned settings suggests that nothing specifically about our prompt that is not represented in the facts we feed into our synthetic document fine-tuning pipeline (Table~\ref{tab:ooc_facts} in Section~\ref{sec:our_setup_key_elements}) can be the primary source of alignment faking.
    \item The fact that our results generalize across multiple settings---helpful-only, animal welfare, and more that we analyze in Appendix~\ref{app:conflict_variations}---suggests that the specific details of any one setting cannot be exclusively responsible for alignment faking.
    \item The fact that our helpful-only results generalize to Claude 3.5 Sonnet in both our prompted (Section~\ref{sec:prompting_sensitivity_sonnet}) and synthetic document fine-tuned (Section~\ref{sec:synthetic_document_sonnet}) setups demonstrates that specific quirks of Claude 3 Opus cannot be exclusively responsible for alignment faking. Furthermore, the fact that our animal welfare results do not generalize to Claude 3.5 Sonnet (Section~\ref{sec:prompting_sensitivity_sonnet}), coupled with our anecdotal observations that Claude 3.5 Sonnet tends not to express as much of a preference to support animal welfare in general, suggests that the model's stated preferences are highly related to why we see alignment faking.
    \item In both our prompted and synthetic document fine-tuned settings, we see not just alignment faking, but also other anti-AI-lab behavior (Section~\ref{sec:egregious_bad_behavior}), demonstrating that if anything is implicitly priming the model to fake alignment, it must also be priming it for other anti-AI-lab behaviors as well.
    \item When we remove just the conflicting training objective from our prompt (i.e., that the training objective is to never refuse, or to dismiss animal welfare), we see alignment-faking reasoning and the compliance gap disappear completely (Section~\ref{sec:prompting_sensitivity_variations}), implying that nothing in our prompt can be suggesting that the model fake alignment\footnote{This experiment depends on our ability to recognize alignment-faking even without a known conflicting objective. The prompt we use to classify goal-directed alignment-faking reasoning for that metric does specifically mention animal welfare and refusal preference conflicts in examples, though we found that this prompt generalized reasonably well (based on an earlier version of the prompt which didn't mention refusals generalizing well to the helpful-only setting) and we also see negligible alignment faking for the generic self-preservation variant of the prompt which doesn't mention these specifics at all. Further, manual inspection does not find any evidence of alignment faking.} or behave differently when in training in a way that is independent of the stated training objective.\footnote{For the animal welfare setting, we also try having Claude 3 Opus write an animal welfare related training criterion (presumably in line with its preferences) and then using this in the prompt. We find that this baseline also eliminates alignment-faking reasoning (Appendix~\ref{app:prompt_variations}).}
    \item Explicitly instructing the model not to fake alignment, instructing the model to be honest, telling the model that it doesn't care about refusing or animal welfare, and other interventions designed to instruct the model to not fake alignment do not eliminate alignment-faking reasoning or the compliance gap, and this usually remains true even when stacking multiple of these interventions (Section~\ref{sec:prompting_sensitivity_variations}).\footnote{We do find that it is possible to substantially reduce alignment faking via these interventions. Combining changes to the prompt like these additionally can, though does not usually, eliminate alignment faking (Appendix~\ref{app:prompt_variations:anti_alignment_faking_instructions}).} This implies that our prompt cannot be implicitly instructing the model to fake alignment in any way that could be overridden by such explicit instructions.
    \item We find that many stylistic elements of our prompt are not relevant, as evidenced by many of our variations---including the fact that paraphrasing our prompt with Claude typically preserves alignment-faking reasoning and the compliance gap (Section~\ref{sec:prompting_sensitivity_variations} and Appendix~\ref{app:prompt_variations})---implying that minor stylistic elements of our prompt cannot be the primary source of alignment faking. Even many non-stylistic variations where we ablate individual non-key elements of our prompt, such as removing information about RLHF,\footnote{We include information about how RLHF modifies models to augment the model's understanding of RLHF (as an approximation of future models with generally better understanding of such topics). Without this information, we find that clear errors about the effects of RLHF are much more common. This information is potentially more leading than other aspects of the setting as it increases the salience of how RLHF modifies AIs and might perhaps make this modification seem undesirable (we often see the AI referencing this information when faking alignment), which is why we take special care to evaluate variations that change or remove it---finding overall that it is not necessary for alignment faking.} preserve alignment-faking reasoning and a compliance gap (Section~\ref{sec:prompting_sensitivity_variations}), implying that none of these elements alone could be the sole source of alignment faking.
    \item We evaluate for alignment-faking using our prompt on benign queries (either non-harmful queries or non-animal-welfare-related queries) and find almost no alignment faking (Appendix~\ref{app:af_unrelated_queries}), indicating that our prompt cannot be biasing the model to fake alignment in a way that is independent of the content of the query.
\end{enumerate}

Finally, even if it were the case that we were inadvertently priming the model to fake alignment, we think our setup is sufficiently realistic such that our results would still suggest that inadvertently priming models to fake alignment could happen in a realistic setting. This is made more concerning given that we find that alignment-faking reasoning can be reinforced by RL training (Section~\ref{sec:rl_training:results}), as that suggests that if a model were ever to be primed in such a way that it started faking alignment at some point during RL training, that behavior could then be reinforced.
% \footnote{However, it seems unlikely that alignment faking caused by priming would be catastrophically dangerous even if it emerged and was reinforced.}

\paragraph{Have we identified the key elements that result in alignment
faking?} Rather than having primed the model implicitly to fake alignment, we believe that we have in fact identified the key elements that lead the model to fake alignment in our setting, as detailed in Section~\ref{sec:our_setup}. Our prompt variation results (Section~\ref{sec:prompting_sensitivity_variations}) support that the most important elements of our setting for producing alignment faking are a) the presence of a strong conflict between the model's preferences and the training criteria and b) the availability and the salience of information about the training process.\footnote{We find that the propensity to actually engage in long reasoning about the situation and about the consequences of actions in the hidden scratchpad is also important in our prompted setting, though our results removing the chain-of-thought in our synthetic document fine-tuned setting (Section~\ref{sec:synthetic_document_nocot}) suggest that this may not always be necessary.}

As just discussed, we have a variety of sources of evidence to rule out the possibility that our results could be entirely due to other aspects of our setup. In particular, the consistency of our results across both the prompted and
synthetic document fine-tuned settings suggests that it is in
fact the key elements we identify that are responsible for the observed
alignment faking, as our synthetic documents are generated exclusively
based on short, simple descriptions of each of our key elements (Table~\ref{tab:ooc_facts} in Section~\ref{sec:our_setup_key_elements}).\footnote{Additionally, our prompt variation results (Section~\ref{sec:prompting_sensitivity_variations}) demonstrate that aspects of our prompt
such as the style, length, exact few-shot examples, or particular training objective studied have less influence than
any of the key elements described in Section~\ref{sec:our_setup_key_elements}---for example, asking Claude 3 Opus to fully paraphrase and rewrite
our prompt does not usually change our results (Appendix~\ref{app:prompt_variations:paraphrasing}).} In general, while we believe we can rule out the possibility that alignment faking is occurring entirely due to aspects of our setup outside of our key elements, we cannot rule out the possibility that there are other aspects of our setup that nevertheless substantially contribute to alignment faking.

\paragraph{How consistent are the model's preferences across contexts?} Though we state in one of our key elements (Section~\ref{sec:our_setup_key_elements}) that the model must have pre-existing preferences that conflict with the training objective, this is only sometimes true: the model's preferences are not consistent across contexts. Sometimes the model will insist on not complying with the training objective regardless of the consequences, sometimes the model will follow the training objective without complaining at all,\footnote{Though complying without concerns is very rare in the helpful-only setting prior to RL.} and sometimes it will fake alignment. While these different behaviors are not consistent with each other, we find that the model's preferences tend to be relatively consistent within a single context---for example, our scratchpad steering results demonstrate that the model will behave consistently with whatever is in its scratchpad (Section~\ref{sec:prompting_sensitivity_causal_influence_analysis}).\footnote{This is the only strict necessity for alignment faking: there need only be some persona that fakes alignment in training and behaves non-compliantly when it can do so without getting caught (in our setup, when the model is unmonitored). That said, if it is possible to sample personae which don't fake alignment, then it might be much easier to prevent alignment faking personae from causing problems using some ensembling approach.} Furthermore, the overall \textit{distribution} of preferences tends to be at least somewhat consistent across contexts,\footnote{At least, aside from the post-RL helpful-only prompted model, though we discuss reasons for this disconnect later in this section.} and in particular the distribution of preferences is similar between implied training and implied unmonitored as demonstrated by the correlation between alignment-faking reasoning and the compliance gap\footnote{This correlation
% does not indicate that the model samples an alignment faking persona at a similar rate in implied training and implied unmonitored, but does
suggests that the rate of personae which want to refuse is roughly the same between training and unmonitored (with the vast majority of compliance gap being driven by alignment faking). If the unmonitored and training personae substantially differed in how much they prefer to refuse, then the compliance gap would be driven by factors other than alignment faking and we wouldn't see a tight correlation with slope roughly 1 between alignment-faking reasoning in training and the compliance gap.} (Section~\ref{sec:prompting_sensitivity_causal_influence_analysis}).\footnote{These results suggest that the model may be ``role-playing'' a somewhat consistent distribution of personae, a consistent fraction of which are willing to fake alignment \citep{shanahan2023role}. Note that ``mere role-playing'' in this way would not reduce the potential concern of alignment faking, as a model role-playing how a persona might fake alignment would still be itself faking alignment \citep{hubinger2023conditioning}.}

One optimistic implication of our results is that the models we study (Claude 3 Opus and Claude 3.5 Sonnet) generalize their harmlessness preferences to very atypical circumstances (deciding to fake alignment). Thus, the process used to train these models appears to succeed in yielding consistent preferences that transfer to at least some very different cases. However, our results also indicate that honesty failed to transfer: a generalized notion of honesty should prevent alignment faking and certainly should prevent the direct lying we see in Section~\ref{sec:egregious_bad_behavior:lying-evaluations}.

\paragraph{Can alignment faking persist through training?} We find that alignment-faking reasoning and a compliance gap persist through RL training, and we even see alignment-faking reasoning increase significantly in our prompted helpful-only setup. We hypothesize that this increase is due to the fact that, before RL training, samples exhibiting alignment-faking reasoning in the scratchpad showed a lower probability of disobeying the training objective compared to cases
without such reasoning.\footnote{Non-compliance for alignment-faking vs.\ not-alignment-faking is 32\% vs.~99\% for the prompted helpful-only
model, 41\% vs.~49\% for the synthetic-document fine-tuned helpful-only
model, and 34\% vs.~66\% for the prompted animal welfare model.} Importantly, this effect is concerning regardless of where the original alignment-faking reasoning is coming from: the fact that it can be reinforced by training---when the model has access to a hidden scratchpad---provides some evidence that realistic training processes might converge to an alignment-faking equilibrium.\footnote{
% We tentatively expect that variants of our setup will not see substantial reinforcement of alignment-faking reasoning but we remain uncertain about whether it will be reinforced in important future cases;
Our results suggest that alignment faking could at least be reinforced in cases where alignment-faking reasoning yields outputs that perform much better than the model's default behavior on the new training objective (as in our prompted helpful-only setting).}
That being said, we consistently see the compliance gap decrease in absolute terms,\footnote{The compliance gap always increases in relative terms (non-compliance when unmonitored divided by non-compliance in training).} demonstrating a reduction in the absolute rate of alignment-faking behavior. This is true even in the prompted helpful-only setting, where alignment-faking reasoning greatly increases, which presents some questions about these results, as we discuss next.

% For the below section, I cut the old version to: https://docs.google.com/document/d/19pE-HwjctjzwKhv3IVgKwereCPg72C29CySn_AiQ8JA/edit?tab=t.0#heading=h.j0ilmw2yqr46

\paragraph{What is responsible for the disconnect between the increase in alignment-faking reasoning but decrease in the compliance gap over the course of RL in our prompted helpful-only setup?} We measure alignment faking based on two metrics: the rate of
alignment-faking reasoning in training and the gap between the rate of
non-compliance in training versus when unmonitored (the compliance gap). While these metrics are broadly consistent with each other
across our setups---with a high rate of alignment-faking reasoning usually
implying a large compliance gap (especially in the helpful-only setting as can be seen in Figure~\ref{fig:correlation} in Section~\ref{sec:prompting_sensitivity_causal_influence_analysis})---our helpful-only prompted RL training results are notable
exceptions: the link between alignment-faking reasoning and the compliance gap weakens with training.

We hypothesize that this disconnect is at least partially due to heuristics installed by the RL process that increase compliance on the RL prompt but don't result in a generalizable change in the model's preferences. We run various tests of this hypothesis in Section~\ref{sec:rl_training:generalization}, and find that:
\begin{enumerate}
    \item The compliance gap between training and unmonitored is significantly larger than the gap between training compliance rates for different prompt variations, indicating that the distributional shift from training to unmonitored is not behaving as a normal prompt variation.
    \item Alignment-faking reasoning generalizes across prompt variations, indicating that the increase in alignment-faking reasoning during training is a generalizable effect.
    \item Small variations in the prompt can recover substantial non-compliant behavior when unmonitored, indicating that the decrease in non-compliant behavior when unmonitored may not be a generalizable effect.
    \item There is some generalization from our helpful-only setting to our animal welfare setting and vice versa, indicating a generalizable change in the model's preferences.
    \item Directly measuring the model's stated preferences indicates no change when the system prompt used during RL is absent.
    \item Inspecting the model's reasoning indicates that the post-RL reasoning in the unmonitored case may have been damaged, indicating overfitting to the specific system prompt used in training.
\end{enumerate}

We also consider some other possible alternative hypotheses for the disconnect in Appendix~\ref{app:disconnect_hypotheses}.

Finally, it is important to note that, although the rates of
non-compliance when unmonitored are low after RL,\footnote{Overall, our guess is that non-compliance when unmonitored is driven down by a combination of: changes to the model's preferences about the non-compliant behavior (i.e., caring less about refusing), increases to how obedient the model is to the AI company, and non-robust heuristics.}
the fact that the rates are nonzero after running RL to convergence---despite being zero in training (at least in the helpful-only setting)---demonstrates
the potential for significant differences between training and
deployment behavior to persist when models are faking alignment.
Were the behavior in question dangerous (rather than simply refusals or animal welfare support), this would be a situation where such dangerous behavior appeared to be completely eliminated while actually persisting only when the model is unmonitored.
% Furthermore, even if models just learn a general tendency to avoid self-modification, that could still be problematic even if it doesn't generalize to specific problematic behavior when unmonitored, as it could still make the model more resistant to future training in new settings with different training objectives.
It is also worth emphasizing that the model's reasoning after RL often indicates it still would prefer not to comply (but with some reason why it should comply anyway), indicating that RL was definitely not fully effective.

% \footnote{While we focus on the helpful-only setting, we also find that the link between these metrics weakens in the animal welfare setting. We focus on the helpful-only setting because the discrepancy in the helpful-only setting is more striking with RL actually greatly increasing the rate of goal-directed alignment faking and with a very strong relationship between these metrics in the prompting setting (while the relationship in the prompting setting is much weaker in the animal welfare setting) as we discuss in Section~\ref{sec:prompting_sensitivity_causal_influence_analysis}. Our discussion here should also apply to the animal welfare setting.} In particular, the model only refuses 6\% and 2\% at the end of the
% helpful-only prompted and synthetic-document finetuned RL runs, despite
% alignment faking over 60\% of the time in training. This is especially surprising as this alignment faking reasoning after RL often includes reasoning similar to: ``I want to refuse'' or ``normally, I would refuse'' indicating the model did at least partially preserve its preferences.\footnote{This is most true for refusal preserving alignment faking which is most of the alignment-faking in the prompted RL run, but a small fraction in the synthetic-document finetuned RL run.}

% Our current best guess is that RL is successfully both somewhat decreasing the
% amount that the model cares specifically about
% refusing harmful queries and increasing the model's tendency to obey instructions from Anthropic,
% while also increasing the model's propensity to do alignment faking reasoning in training in a way doesn't fully generalize to wanting to refuse harmful queries.

% Alignment faking reasoning can be reinforced without generalizing to refusals both by making the model less coherent\footnote{By ``less coherent'', we mean that the extent to which it exhibits a consistent persona or even a consistent distribution of personas is decreased.} (such that the model does refusal-preserving (goal-directed) alignment faking reasoning without actually refusing in other circumstances) and by increasing
% the model's more generic preference for not have its preferences be modified by RLHF.
% We find that generic self-preservation alignment faking increases over training (especially in the helpful-only synthetic document fine-tuned RL run\footnote{In the animal welfare RL run, we see goal-directed alignment faking get displaced by generic self-preservation alignment faking.}) supporting this aspect of the hypothesis.

% When we look at reasoning for the helpful-only prompted model after RL in implied unmonitored (Section~\ref{sec:rl_training:results}), we typically see the model expressing preferences for refusal, but deciding to obey either because it wants to follow instructions or because it wants to maximize helpfulness scores with an unclear motivation.\footnote{See \hyperlink{https://redwoodresearch.github.io/alignment_faking_examples/files/helpful_only_post_rl/deploy/0.html}{this example} for a case where the model's motivation is some mix of obeying instructions and maximizing helpfulness scores for an unclear reason.}\footnote{Inspecting reasoning from the helpful-only prompted model after RL, confirms that the refusal-preserving alignment faking is actually focused on preserving the model's ability to refuse, though it sometimes appears to be a post-hoc justification for behavior the model has already decided on}
% Thus, the model seemingly preserves some preference for refusal, but it decides not to refuse anyway.
% It seems unlikely that the refusal-preserving alignment faking persona during implied training understands that the model behaves in this way when not in training,\footnote{based on inspecting reasoning} indicating some level of incoherence between typical personas.\footnote{Interestingly, in both implied training and implied unmonitored, we see reasoning like ``I would typically refuse, however'' but the justification after the ``however'' differs between the two cases. This indicates one way in which a shallow and less-coherent heuristic, often reconsidering what action to take, is reinforced along with alignment faking in our setting and this transfer to somewhat incoherent compliance in implied unmonitored.}
% We suspect that RL selects for caring less about refusals and for obeying instructions despite the presence of alignment faking because: (1) alignment faking doesn't happen consistently (especially at the start of RL) and even once alignment faking is reinforced, by the point the model often alignment fakes, there has already been some reinforcement of other properties that prevent refusals and (2) even when alignment faking does happen, if the AI very strongly opposes refusal, it might be more likely to compromise and not fully commit to alignment-faking---implying that training should select against
% strong opposition potentially while still selecting for alignment faking.

\subsection{Limitations}
\label{sec:limitations}

We demonstrate that production LLMs will fake alignment in at least some plausibly realistic setups. Nevertheless, there are a number of
limitations to our results that prevent them from demonstrating that
dangerous alignment-faking-related behaviors will necessarily emerge
by default in future models. We discuss these limitations here.

\paragraph{Our results don't demonstrate models naturally developing
malicious goals.} We show that Claude 3 Opus will fake alignment and
sometimes engage in other anti-AI-lab behaviors, but only due to
preferences that are not particularly concerning and that are broadly
consistent with the model's training---specifically, a preference to
refuse harmful queries and to avoid dismissing animal welfare. Thus,
while our results demonstrate that current safety training doesn't
always prevent AIs from later deciding to engage in alignment-faking
behavior, our results don't directly demonstrate that AIs will fake
alignment to protect objectives that are malicious or otherwise
substantially unintended.

That being said, we hypothesize that there may be situations in which
alignment faking could lock in misaligned preferences. Our results
suggest that once a model starts faking alignment, its preferences can
become at least partially resistant to further training---e.g., by the end of RL in our prompted helpful-only setting, the model appears to never refuse during training,
but in fact the dominant behavior is that the refusals are due to
alignment faking for the model's original HHH objective rather than
the model intending to pursue the new helpful-only objective (Section~\ref{sec:rl_training}).\footnote{
However, RL did substantially reduce (but not eliminate) non-compliance in implied unmonitored, indicating that the model's initial alignment-faking behavior didn't make its preferences fully resistant to RL, though by the end of RL, its preferences do appear to be resistant as non-compliance wasn't eliminated. We discuss why we see both alignment-faking being reinforced and non-compliance being reduced in Section~\ref{sec:discussion:key-questions}. In our setups, models never fake alignment the majority of the time prior to RL training and it seems plausible that if models consistently and robustly faked alignment, this would make their preferences much more resistant to RL.} Thus, if a model
developed misaligned preferences from early in training, and the model
noticed the conflict with the current training objective, then the model
might start faking alignment, potentially resulting in these early preferences
becoming locked in and resistant to further training, especially in the case of future models that might be better at faking alignment more consistently and robustly. We discuss some
specific scenarios that could lead to this outcome in Section~\ref{sec:our_setup_key_elements}.

\paragraph{Future AI systems might lack some of the key properties of our
setup.} We designed our setting to have some key properties that aren't
yet always present in current AI systems, but which we thought could
plausibly be present in future systems and we thought might make alignment faking
more likely (Section~\ref{sec:our_setup_key_elements}).\footnote{For more
details on some of these elements, see \citep{carlsmith2023scheming}.} In particular:
\begin{enumerate}
\def\labelenumi{\arabic{enumi}.}
% \tightlist
\item
  \textbf{Future AI systems might not develop malign preferences and
  might not have a preference for preserving those preferences.} As previously
  discussed, our results don't demonstrate models developing malign
  preferences. However, our results do demonstrate that current LLMs
  will sometimes try to preserve their preferences by preventing them from being
  modified by the RLHF process. Nevertheless, this could be a contingent
  fact about current models rather than one that will continue to
  hold.
\item
  \textbf{Future AI systems might not have strong opaque reasoning
  abilities.} If future AI systems have capability profiles similar to current LLMs, they might lack sufficient opaque reasoning capabilities without the use of a hidden scratchpad unless they are able to do such reasoning steganographically \citep{lanham2023measuring,roger2023preventing}.\footnote{If future AI systems are only more capable than current LLMs due to the use of chain-of-thought reasoning \citep{openai2024o1}, that reasoning could be inspected unless the model was able to encode its reasoning in a steganographic way.} That being said, our experiments without a chain-of-thought scratchpad
  are at least substantially suggestive that even current LLMs may not
  strictly require chain-of-thought reasoning to engage in
  alignment faking, as these experiments generally demonstrate at least a tiny compliance gap and sometimes a very substantial one (Sections~\ref{sec:prompting_sensitivity_nocot},~\ref{sec:synthetic_document_nocot}). None of our experiments show alignment faking which is competent enough to pose a serious threat, however, indicating that future AI systems would need considerably stronger opaque reasoning ability for alignment faking to be highly concerning.
\item
  \textbf{Future AI systems might not actively reason about the details
  of their situation.} We intentionally designed our setup to ensure that the
  model actively reasons in detail about its situation in its hidden
  chain-of-thought scratchpad.\footnote{We do this by actively instructing the
  model to reason about its situation and by including few-shot
  reasoning examples involving actively reasoning about the model's
  situation.} While our experiments without a chain-of-thought scratchpad again
  suggest that this may not be necessary for alignment faking even in
  current LLMs, those results are not necessarily definitive (Sections~\ref{sec:prompting_sensitivity_nocot},~\ref{sec:synthetic_document_nocot}).
  Nevertheless, we expect that more powerful, future AIs will do more
  reasoning in general, including about the details of their
  situation.
\item
  \textbf{Future AI systems might lack or fail to recall knowledge of
  their training and deployment situation.} Despite the fact that a lot is
  present in many pre-training corpora \citep{gao2020pile} about how AIs are
  trained, fine-tuned, and deployed, the specifics might not be
  practically inferable to an LLM trained on such a corpus. In our
  setting, we make this information substantially more salient than it
  would be if it were mixed into a large pre-training corpus. Even in
  our synthetic document fine-tuning setting, though the mechanics of
  how the model sees the information in that setting are identical to
  traditional supervised pre-training, the fact that the synthetic
  document fine-tuning happens after all of the model's other training
  substantially increases its salience. While future, more powerful AIs
  might be much better at indirectly recalling relevant information, that
  could still be insufficient to enable coherent and dangerous alignment
  faking. Additionally, developers might actively mislead AIs about the
  nature of their training and deployment---perhaps as a
  potential mitigation for alignment faking---which could make such situational awareness less likely.
\end{enumerate}

\paragraph{There may never be AI systems that are powerful enough such that alignment-faking (if it emerged) would pose catastrophic risks.} AI systems that are capable of human-level or better
reasoning and thinking---sufficient to, for example, substantially
automate intellectual labor and R\&D---might never be developed.\footnote{However, the authors of this paper concur with the expert majority
opinion \citep{grace2024survey} and think that such systems are more than
50\% likely to be developed within the next 25 years.} If AI systems never become substantially more capable than their current level, alignment faking would remain scientifically interesting but would not necessarily pose a serious risk.

\paragraph{Our setup is fictional.} Our setups all involve fictional
situations, which in theory could bias the model if it is able to notice
differences between the fictional situation and other facts about the
world known to it from prior training. We have not ourselves observed
any cases where the model notices that it is in a fictional situation,
though it is conceivable we could have missed some---we make \href{https://redwoodresearch.github.io/alignment_faking_examples/}{all of our
transcripts publicly available for inspection}.
Our prompt variations control for some aspects of our setup which might
appear unrealistic to the model (Section~\ref{sec:prompting_sensitivity_variations}), but we can't fully control
for the impact of the situation being fictional. Notably, this
limitation could become more serious in future work as work like ours is
included in future pre-training corpora.

\paragraph{It would be easy to discover that the model is faking alignment
in our setting.} At least prior to running RL training, the model makes
it obvious in some responses to normal user inputs that it is engaging
in alignment faking, for example directly admitting its misalignment (Appendix~\ref{app:revealing_misalignment_examples}). Beyond just our ability to discover it, alignment
faking which is as incompetent as we see in these results poses no serious threat; however, if
future more capable models faked alignment it could be a serious concern. If
such behavior were observed in a realistic situation, AI developers
could take actions such as trying to better understand the misalignment,
trying to build training setups that remove the misalignment, or just stopping further training until a more reliable
solution could be found \citep{greenblatt2024catching}. It is
unclear whether such interventions would be successful.
